%\input{\fullpath delimiters}
%\input{\fullpath summations}

% other definitions
% stray commands



\endinput  %  -  -  -  -  -  -  -  -  -  -  -  -  -  -  -  -  -  -  -  -

%\tiny
%\scriptsize
%\footnotesize
%\small
%\normalsize
%\large
%\Large
%\LARGE
%\huge
%\Huge
%% shades of gray
\definecolor{darkgray}{gray}{0.25}
\definecolor{medgray} {gray}{0.75}
\definecolor{ltgray}  {gray}{0.90}
\definecolor{vltgray} {gray}{0.95}

%\definecolor{rangecolor} {blue}
%\definecolor{nullcolor}  {red}

%% basic color directives
\newcommand{ \bl }[1]        {{\color{blue}   {#1}}}
\newcommand{ \bk }[1]        {{\color{black}  {#1}}}
\newcommand{ \rd }[1]        {{\color{red}    {#1}}}
\newcommand{ \gr }[1]        {{\color{green}  {#1}}}
\newcommand{ \pr }[1]        {{\color{purple} {#1}}}
\newcommand{ \mg }[1]        {{\color{medgray}{#1}}}

%% numbers blue
\newcommand{ \bzero }[0]     { \bl{ 0 } }
\newcommand{ \bone }[0]      { \bl{ 1 } }
\newcommand{ \btwo }[0]      { \bl{ 2 } }
\newcommand{ \bthree }[0]    { \bl{ 3 } }
\newcommand{ \bfour }[0]     { \bl{ 4 } }
\newcommand{ \bfive }[0]     { \bl{ 5 } }
\newcommand{ \bsix }[0]      { \bl{ 6 } }
\newcommand{ \bseven }[0]    { \bl{ 7 } }
\newcommand{ \beight }[0]    { \bl{ 8 } }
\newcommand{ \bnine }[0]     { \bl{ 9 } }
\newcommand{ \bminus }[0]    { \bl{ - } }
\newcommand{ \bstar }[0]     { \bl{ * } }
\newcommand{ \bi }[0]        { \bl{ i } }
\newcommand{ \bmone }[0]     { \bl{-1 } }
\newcommand{ \bmi }[0]       { \bl{-i } }
\newcommand{ \bmo }[0]       { \bl{-1 } }

\newcommand{ \bdots }[0]     { \bl{ \dots } }


%% numbers red
\newcommand{ \rzero }[0]     { \rd{ 0 } }
\newcommand{ \rone }[0]      { \rd{ 1 } }
\newcommand{ \rtwo }[0]      { \rd{ 2 } }
\newcommand{ \rthree }[0]    { \rd{ 3 } }
\newcommand{ \rfour }[0]     { \rd{ 4 } }
\newcommand{ \rfive }[0]     { \rd{ 5 } }
\newcommand{ \rsix }[0]      { \rd{ 6 } }
\newcommand{ \rseven }[0]    { \rd{ 7 } }
\newcommand{ \reight }[0]    { \rd{ 8 } }
\newcommand{ \rnine }[0]     { \rd{ 9 } }
\newcommand{ \rminus }[0]    { \rd{ - } }
\newcommand{ \rstar }[0]     { \rd{ * } }
\newcommand{ \rmone }[0]     { \rd{-1 } }

\newcommand{ \rdots }[0]     { \rd{ \dots } }

%% numbers medium gray
\newcommand{ \gzero }[0]     { \mg{ 0 } }
\newcommand{ \gone }[0]      { \mg{ 1 } }
\newcommand{ \gtwo }[0]      { \mg{ 2 } }
\newcommand{ \gthree }[0]    { \mg{ 3 } }
\newcommand{ \gfour }[0]     { \mg{ 4 } }
\newcommand{ \gfive }[0]     { \mg{ 5 } }
\newcommand{ \gsix }[0]      { \mg{ 6 } }
\newcommand{ \gseven }[0]    { \mg{ 7 } }
\newcommand{ \geight }[0]    { \mg{ 8 } }
\newcommand{ \gnine }[0]     { \mg{ 9 } }
\newcommand{ \gminus }[0]    { \mg{ - } }
\newcommand{ \goplus }[0]    { \mg{ \oplus } }

%% numbers black
\newcommand{ \bs }[0]        { {\bk{*}} }
\newcommand{ \bkzero }[0]    { {\bk{0}} }

% nums
\newcommand{ \bnum }[0]      { \bl{ \num } }
\newcommand{ \rnum }[0]      { \rd{ \num } }
\newcommand{ \gnum }[0]      { \mg{ \num } }

\endinput  %  -  -  -  -  -  -  -  -  -  -  -  -  -  -  -  -  -  -  -  -
% programming environment
\newcommand{\scru}[1]     {\bl   {\texttt{#1}}}
\newcommand{\scrk}[1]     {\bk   {\texttt{#1}}}
\newcommand{\scrg}[1]     {\gr   {\texttt{#1}}}
\newcommand{\scrr}[1]     {\rd   {\texttt{#1}}}
\newcommand{\scrv}[1]     {\dg   {\texttt{#1}}}
\newcommand{\scrz}[1]     {\mg   {\texttt{#1}}}
\newcommand{\scrc}[1]     {\textcolor{cyan}  {\texttt{#1}}}
\newcommand{\scrp}[1]     {\textcolor{purple}{\texttt{#1}}}
\newcommand{\scry}[1]     {\textcolor{yellow}{\texttt{#1}}}
\newcommand{\scrw}[1]     {\textcolor{white} {\texttt{#1}}}

%
\newcommand{\mySize}[1]   {{\normalsize{#1}}}
\newcommand{\myCanvas}[1] {{\setbeamercolor{background canvas}{bg=black}{#1}}}
\newcommand{\myRasa}[1]   {\myCanvas{\mySize{#1}}}
\newcommand{\myFrame}[2]  {\myRasa{\begin{frame} \frametitle{#1} #2 \end{frame}}}
\newcommand{\myFoo}[2]    {\myFrame{#1}  \program \ \\[10pt] #2 \eprogram}


% program environment
\newcommand{\program}[0]  {\scrz{program myProgram}}
\newcommand{\eprogram}[0] {\scrz{end program myProgram}}
\newcommand{\inone}[0]    {\scrz{implicit none}}

% tabs
%\newcommand{\tab}[0]        {\phantom{mm}}
\newcommand{\tab}[0]      {\ \ }
\newcommand{\spacer}[0]   {\ps\text{::}\ps}
\newcommand{\ps}[0]       {\phantom{m}}

% is equal?
\newcommand{\iseq}[0]    { \overset{?}{=} }


% spacings and such
\newcommand{\wfour}[0]   {1.05in}


\endinput  %  -  -  -  -  -  -  -  -  -  -  -  -  -  -  -  -  -  -  -  -

% https://engineering.purdue.edu/ECN/Support/KB/Docs/LaTeXChangingTheFont
% \tiny
% \scriptsize
% \footnotesize
% \small
% \normalsize
% \large
% \Large
% \LARGE
% \huge
% \Huge% programming enviro
\newcommand{\paren}[1]   { \left(  #1 \right) }
\newcommand{\inner}[2]   { \langle #1 \overline{#2} \rangle }
\newcommand{\innerx}[2]  { \langle #1 | #2 \rangle }
\newcommand{\brac}[1]    { \left[  #1 \right] }
\newcommand{\lst}[1]     { \left\{ #1 \right\} }

\endinput  %  -  -  - -  -  - -  -  - -  -  - -  -  - -  -  - -  -  - -  -  -
\input{\pathname/"matrix basics"}
\newcommand{\abs}[1]     {\left| #1 \right|}
\newcommand{\norm}[1]    {\left\lVert #1 \right\rVert}
\newcommand{\normo}[1]   {\left\lVert #1 \right\rVert_{1}}
\newcommand{\normt}[1]   {\left\lVert #1 \right\rVert_{2}}
\newcommand{\normi}[1]   {\left\lVert #1 \right\rVert_\infty}
\newcommand{\normp}[1]   {\left\lVert #1 \right\rVert_{p}}
\newcommand{\normf}[1]   {\left\lVert #1 \right\rVert_F}
\newcommand{\normtL}[1]  {\left\lVert #1 \right\rVert_{L^{2}}}
\newcommand{\normtl}[1]  {\left\lVert #1 \right\rVert_{l^{2}}}

\newcommand{\norms}[1]   {\left\lVert #1 \right\rVert^{2}}
\newcommand{\normos}[1]  {\left\lVert #1 \right\rVert_{1}^{2}}
\newcommand{\normts}[1]  {\left\lVert #1 \right\rVert_{2}^{2}}
\newcommand{\normis}[1]  {\left\lVert #1 \right\rVert_\infty^{2}}
\newcommand{\normps}[1]  {\left\lVert #1 \right\rVert_{p}^{2}}
\newcommand{\normfs}[1]  {\left\lVert #1 \right\rVert_F^{2}}
\newcommand{\normtLs}[1] {\left\lVert #1 \right\rVert_{L^{2}}^{2}}
\newcommand{\normtls}[1] {\left\lVert #1 \right\rVert_{l^{2}}^{2}}


\endinput %-------------------------------
% bys
\newcommand{\by}[2]      {#1 \times #2}
\newcommand{\byy}[1]     {#1 \times #1}
\newcommand{\bymn}[0]    {\by{m}{n}}
\newcommand{\bymm}[0]    {\byy{m}}
\newcommand{\bynn}[0]    {\byy{n}}
\newcommand{\bynm}[0]    {\by{n}{m}}
\newcommand{\bymr}[0]    {\by{m}{\rho}}
\newcommand{\byrn}[0]    {\by{\rho}{n}}

% vector spaces
\newcommand{\real}[1]    {\mathbb{R}^{#1}}
\newcommand{\cmplx}[1]   {\mathbb{C}^{#1}}
\newcommand{\either}[1]  {\cmplx{#1}}
\newcommand{\ir}[0]      {\in\real{}}
\newcommand{\ic}[0]      {\in\cmplx{}}
\newcommand{\icm}[0]     {\in\cmplxm}
\newcommand{\icn}[0]     {\in\cmplxn}
\newcommand{\icmn}[0]    {\in\cmplxmn}
\newcommand{\irmn}[0]    {\in\realmn}
\newcommand{\icmnr}[0]   {\in\cmplxmnr}
\newcommand{\ints}[0]    {\mathbb{Z}}
\newcommand{\natnum}[0]  {\mathbb{N}}

\newcommand{\iints}[0]   {\in \mathbb{Z}}
\newcommand{\inatnum}[0] {\in \mathbb{N}}

\newcommand{\realall}[3] {\real{\by{#1}{#2}}_{#3} }
\newcommand{\cmplxall}[3]{\cmplx{\by{#1}{#2}}_{#3} }

\newcommand{\realn}[0]   {\real{n}}
\newcommand{\realm}[0]   {\real{m}}
\newcommand{\realmn}[0]  {\real{\bymn}}
\newcommand{\realnn}[0]  {\real{\byy{n}}}
\newcommand{\realmm}[0]  {\real{\byy{m}}}
\newcommand{\realmmr}[0] {\real{\byy{m}}_{\rho}}
\newcommand{\realmmm}[0] {\realmm_{m}}

\newcommand{\cmplxn}[0]  {\cmplx{n}}
\newcommand{\cmplxm}[0]  {\cmplx{m}}
\newcommand{\cmplxnn}[0] {\cmplx{\byy{n}}}
\newcommand{\cmplxmm}[0] {\cmplx{\byy{m}}}
\newcommand{\cmplxmn}[0] {\cmplx{\bymn}}
\newcommand{\cmplxmr}[0] {\cmplx{\bymr}}
\newcommand{\cmplxrn}[0] {\cmplx{\byrn}}
\newcommand{\cmplxmnr}[0]{\cmplx{\bymn}_{\rho}}
\newcommand{\cmplxmmm}[0]{\cmplxmm_{m}}
\newcommand{\cmplxmnm}[0]{\cmplxmn_{m}}
\newcommand{\cmplxnnr}[0]{\cmplxnn_{\rho}}
\newcommand{\cmplxmnn}[0]{\cmplxmn_{n}}

% spans
\newcommand{\spn}[1]     {\text{sp\,} \lst{ #1 }}


\endinput %-------------------------------

% stray commands
\newcommand{\authorTopa}[0]  {\href{mailto:Daniel.M.Topa@erdc.dren.mil}{Daniel Topa}}
\newcommand{\authorEmbid}[0] {\href{mailto:pfembid@math.unm.edu}{Pedro Embid}}
\newcommand{\engility}[0]    {Engility Corp.\\USACE Engineer Research and Development Center\\Vicksburg MS}

% http://tex.stackexchange.com/questions/34921/how-to-overlap-images-in-a-beamer-slide
\def\Put(#1,#2)#3{\leavevmode\makebox(0,0){\put(#1,#2){#3}}}


\newenvironment{myfancyblock}%
{\begin{center}\begin{tikzpicture}}%
{\end{tikzpicture}\end{center}}%

\newcommand{\opaqueblock}[4]{
\node<#1>[#2=#3] (X) {#4};
}

\newcommand{\onedot}[0]  
{  
  \begin{textblock*}{1cm}(1.136\textwidth,8.9cm)%
   {\color{medgray}{.}}
  \end{textblock*} 
}

\newcommand{\twodots}[0]  
{  
  \begin{textblock*}{1cm}(1.136\textwidth,8.9cm)%
   {\color{medgray}{..}}
  \end{textblock*} 
}

\newcommand{\donedot}[0]  
{  
  \begin{textblock*}{1cm}(1.136\textwidth,8.9cm)%
   {\color{darkgray}{.}}
  \end{textblock*} 
}

\newcommand{\dtwodots}[0]  
{  
  \begin{textblock*}{1cm}(1.136\textwidth,8.9cm)%
   {\color{darkgray}{..}}
  \end{textblock*} 
}

% domains  --   --   --   --   --   --   --   --   --   --   --   --
\newcommand{\bigL}[0]    { L^{2} }
\newcommand{\littlel}[0] { l^{2} }

\newcommand{\innerxL}[2] { \langle #1 | #2 \rangle_{\bigL} }
\newcommand{\innerxl}[2] { \langle #1 | #2 \rangle_{\littlel} }

% action over domains
\newcommand{\domL}[0]    { \Omega }
\newcommand{\doml}[0]    { \sigma }

\newcommand{\intcont}[0] { \Omega = \lst{x\in\real{}\colon -1 \le x \le 1} }
\newcommand{\intdisc}[0] { \sigma = \lst{x\colon -1 \le x \le 1} }

\newcommand{\spaceL}[0]  { \bigL (\Omega) }
\newcommand{\spacel}[0]  { \littlel \paren{\sigma} }

\newcommand{\intL}[0]    { \bigL[-1,1]}
\newcommand{\intl}[0]    { \littlel[-1,1] }

\newcommand{\intdom}[1]  { \int_{\domL} #1(x) #1^{*}(x) dx }
\newcommand{\sumdom}[1]  { \sum_{x\in \doml} #1(x) #1^{*}(x) \Delta }

\newcommand{\intdomx}[2] { \int_{\domL} #1(x) #2^{*}(x) dx }
\newcommand{\intdomy}[2] { \bl{\int_{\domL}} #1(x) #2^{*}(x) dx }

\newcommand{\sumdomx}[2] { \sum_{x\in \doml} #1(x) #2^{*}(x) \Delta }
\newcommand{\sumdomy}[2] { \sum_{x\in \doml} #1(x) #2^{*}(x) \Delta }
\newcommand{\sumdomz}[2] { \bl{\sum_{x\in \doml}} #1(x) #2^{*}(x)\Delta }

%% Title page
\newcommand{\tpage}[0] 
{
\begin{frame}
  \begin{center}
    \bl{{\LARGE{Orthogonality and Computation}}} \\[10pt]
    {\href{mailto:daniel.m.topa@erdc.dren.mil}{\large{Daniel Topa$^{1,2,3}$ and \href{http://math.unm.edu/people/academic-personnel/pedro-embid}{Pedro Embid}}$^{1}$}} \\[10pt]
    %
    \mg{{\footnotesize{$^{1}$ \href{http://www.unm.edu}{University of New Mexico}, \\\href{http://math.unm.edu}{Department of Mathematics and Statistics}}}} \\
    \mg{{\footnotesize{$^{2}$ \href{https://www.lanl.gov}{Los Alamos National Laboratory}}}} \\[3pt]
    \mg{{\footnotesize{$^{3}$ \href{http://www.engilitycorp.com}{Engility Corporation}, \\\href{http://www.usace.army.mil}{USACE} \href{http://www.erdc.usace.army.mil}{Engineer Research and Development Center},\\\href{http://www.erdc.usace.army.mil/Locations/InformationTechnologyLaboratory.aspx}{Information Technology Laboratory}}}} \\[10pt]
    %
    {\normalsize{\href{http://www.world-academy-of-science.org/worldcomp15/ws/conferences/csc15}{CSC'15\\13th International Conference on Scientific Computing}}} \\[5pt]
    {\normalsize{29 July 2015}}
    \Put(-200,200){\href{http://www.world-academy-of-science.org/worldcomp15/ws/conferences/csc15}{\includegraphics[ height = 2cm ]{../graphics/"logo csc"}}}
  \end{center}
\end{frame}
}

% https://tex.stackexchange.com/questions/6370/how-to-center-a-beamercolorbox
\newcommand{\courtney}[0]
{
\begin{frame}      %  %  %   FRAME   %  %  %
\frametitle{Overview}
  %
  \begin{center}
    \includegraphics[ width = 6.75cm ]{../graphics/caution45.png}
  \end{center}
  %
  \begin{textblock*}{5cm}(3.90cm,5.65cm)%
    \begin{center}
     \LARGE{\bf{{\color{\clr}{CAUTION\\[-2pt]}}}}
     \Large{\bf{{\color{\clr}{when using\\[-3pt]orthogonal functions\\[-3pt]in discrete spaces}}}}
    \end{center}
  \end{textblock*}
  %
  %
\end{frame}
}

\newcommand{\courtneyX}[0]
{
\begin{frame}      %  %  %   FRAME   %  %  %
\frametitle{Overview}
  %
  \begin{center}
    {\transparent{0.2}\includegraphics[ width = 6.75cm ]{../graphics/caution.png}}
  \end{center}
  %
  \begin{textblock*}{5cm}(3.90cm,5.65cm)%
    \begin{center}
     \LARGE{\bf{{\color{\clr}{CAUTION\\[-2pt]}}}}
     \Large{\bf{{\color{\clr}{when using\\[-3pt]orthogonal functions\\[-3pt]in discrete spaces}}}}
    \end{center}
  \end{textblock*}
  %
  %
\end{frame}
}


\newcommand{\final}[0]
{
{\setbeamercolor{background canvas}{bg=black}
\color{darkgray}
\begin{frame}      %  %  %   FRAME   %  %  %
\gr{Linearly independent system}
\frametitle{Off-diagonal entries persist in $l^{2}$}
  %
\begin{equation*}
  %\begin{split}
    \mat{cccc}{
    \dg{\innerx{g_{0}}{g_{0}}} & \gr{\innerx{g_{0}}{g_{1}}} & \cdots & \gr{\innerx{g_{0}}{g_{d}}} \\
    \gr{\innerx{g_{1}}{g_{0}}} & \dg{\innerx{g_{1}}{g_{1}}} & \cdots & \gr{\innerx{g_{1}}{g_{d}}} \\
    \vdots & \vdots & \ddots & \vdots \\
    \gr{\innerx{g_{d}}{g_{0}}} & \gr{\innerx{g_{d}}{g_{1}}} & \cdots & \dg{\innerx{g_{d}}{g_{d}}}
    }
    \mat{c}{ b_{0} \\ b_{1} \\ \vdots \\ b_{d} }
    =
    \mat{c}{ \innerx{f}{g_{0}} \\ \innerx{f}{g_{1}} \\ \vdots \\ \innerx{f}{g_{d}} } \quad \gr{\angle}
    \label{eq:l:full}    
  %\end{split}
\end{equation*}

\rd{Orthogonal system}
\begin{equation*}
  %\begin{split}
    \mat{cccc}{
    \dg{\innerx{g_{0}}{g_{0}}} & \rzero & \cdots & \rzero \\
    \rzero & \dg{\innerx{g_{1}}{g_{1}}} & \cdots & \rzero \\
    \vdots & \vdots & \ddots & \vdots \\
    \rzero & \rzero & \cdots & \dg{\innerx{g_{d}}{g_{d}}}
    }
    \mat{c}{ b_{0} \\ b_{1} \\ \vdots \\ b_{d} }
    =
    \mat{c}{ \innerx{f}{g_{0}} \\ \innerx{f}{g_{1}} \\ \vdots \\ \innerx{f}{g_{d}} } \quad \rd{\perp}
    \label{eq:l:full}    
  %\end{split}
\end{equation*}
  %
  \dtwodots
\end{frame}
}
}

\endinput  %  -  -  -  -  -  -  -  -  -  -  -  -  -  -  -  -  -  -  -  -

%\tiny
%\scriptsize
%\footnotesize
%\small
%\normalsize
%\large
%\Large
%\LARGE
%\huge
%\Huge